\documentclass[11pt]{article}
\usepackage[margin=1in]{geometry}
\usepackage{amsmath,amssymb}
\usepackage{booktabs}
\usepackage{siunitx}
\usepackage{hyperref}
\title{Replication Report: \\ \emph{Ab-initio simulation and experimental validation of beta-titanium alloys} \\ (Raabe, Sander, Fri\'ak, Ma, Neugebauer; arXiv:0811.0157, 2008)}
\author{Automated replication (OpenClaw subagent)}
\date{2026-07-19}
\begin{document}
\maketitle

\begin{abstract}
We replicate the machine-checkable computational claims of Raabe et al. (2008), a DFT+experiment
progress report on low-stiffness biomedical $\beta$-Ti alloys. The paper is \textbf{off-theme}
for the ``texture-spin'' batch it was filed under (``texture'' is one characterized property, not
the subject). Its true ab-initio engine (plane-wave DFT total energies, GGA-PBE96) is HPC-scale
and out of scope; we therefore isolate and reproduce the \emph{closed-form} thermodynamics
(ideal-mixing entropy, finite-$T$ free energy) and the quantitative data/arithmetic claims.
Of four selected claims: C1 (entropy) PASS, C3 (elastic data) PASS (3/3), C4 (composition
conversion) PASS, C2 (finite-$T$ $\beta$-stabilization threshold) PARTIAL. Overall
\textbf{PARTIAL REPLICATION}; Coverage 6/10, Agreement 8/10.
\end{abstract}

\section{Paper claims and method}
The paper combines DFT (GGA-PBE96) with casting/ultrasonic experiments to design $\beta$-Ti alloys
with a low elastic modulus (approaching cortical bone, $\sim$20\,GPa, vs Ti-6Al-4V $\sim$105--110\,GPa).
Key relations:
\begin{align}
E_f^\sigma(\mathrm{Ti}_xX_{1-x}) &= E_\mathrm{tot}^\sigma/N - x\,\mu_\mathrm{Ti}^{hcp} - (1-x)\mu_X^{bcc}, & \text{(1)}\\
S_\mathrm{config}(x) &= -k_B\big[x\ln x + (1-x)\ln(1-x)\big], & \text{(2)}\\
F_f(x,T) &= \langle E_f\rangle - T\,S_\mathrm{config}, & \text{(3)}\\
\sigma &= (\delta E/\delta\varepsilon)/(A\,a_0). & \text{(4)}
\end{align}
(Eq.~2 is printed in the paper without the leading minus; we use the physically-correct positive
form, verified below.)

\section{Selected machine-checkable claims}
\begin{description}
\item[C1] Ideal-mixing configurational entropy (Eq.~2), with maximum $S(0.5)=k_B\ln 2$.
\item[C2] Finite-$T$ free energy (Eq.~3) lowers the $\beta$-stabilization threshold: Ti-Nb from
  $\sim$93 at.\% (T=0) toward $\sim$25 at.\% near the transition temperature; Ti-Mo from
  $\sim$25 at.\% to $\sim$14 at.\%.
\item[C3] Reported experimental Young's moduli: minimum is Ti-30at.\%Nb (72.1\,GPa); $\sim$37\%
  drop vs hcp-Ti reference (114.7\,GPa); trends.
\item[C4] Table-1 wt\%$\leftrightarrow$at\% conversions for all 8 alloys.
\end{description}

\section{Results}

\subsection*{C1 --- Configurational entropy (PASS)}
Our implementation of Eq.~(2) gives $S(0.5)=0.69315\,k_B$, matching $k_B\ln 2=0.69315\,k_B$ to
machine precision; $T\!\cdot\!S$ at the transition temperature ($\sim$1155\,K) is $\sim$69\,meV/atom
at $x=0.5$. The entropy machinery is exact.

\subsection*{C2 --- Finite-$T$ $\beta$-stabilization threshold (PARTIAL)}
We lack the DFT $\langle E_f(x)\rangle$ curve, so we anchor a labelled analytic surrogate to the
paper's stated $T=0$ roots and apply Eq.~(3):
\begin{center}
\begin{tabular}{lccc}
\toprule
System & $T=0$ root (anchor) & finite-$T$ (this work) & finite-$T$ (paper) \\
\midrule
Ti-Nb & 93 at.\% & 79.0 at.\% & $\sim$25 at.\% \\
Ti-Mo & 25 at.\% & 16.6 at.\% & $\sim$14 at.\% \\
\bottomrule
\end{tabular}
\end{center}
Ti-Mo reproduces well; Ti-Nb does not. The gap is a limitation of the symmetric-parabola energy
surrogate (the real Ti-Nb $E_f(x)$ is skewed toward dilute compositions), \emph{not} of Eq.~(3).
Reported honestly, not tuned.

\subsection*{C3 --- Elastic modulus data (PASS, 3/3)}
Quoted experimental moduli (GPa): Ti-Nb \{10:91.2, 20:75.8, 30:72.1\}, Ti-Mo \{10:88.4, 20:106.9\},
hcp ref 114.7.
(a) Minimum measured modulus $=72.1$\,GPa $=$ Ti-30Nb (PASS).
(b) Drop vs hcp ref $=100(114.7-72.1)/114.7=37.1\%$, paper $\sim$37\% (PASS).
(c) Ti-Mo modulus increases with Mo (slope $+1.85$\,GPa/at.\%, $R^2=1.000$, PASS); Ti-Nb
experimental series \emph{decreases} (slope $-0.96$, $R^2=0.889$), consistent with the paper's own
$\alpha/\omega$-contamination caveat for low-solute samples.

\subsection*{C4 --- Composition conversion (PASS)}
Recomputing at.\% from the tabulated wt.\% (standard atomic masses) reproduces the paper's at.\%
column: 18/20 element rows agree to $<0.25$ at.\%; the single 1.02 at.\% mismatch (Ti-20Nb) is a
rounding artifact of the paper's own rounded wt.\% input.

\section{Verdict}
\textbf{PARTIAL REPLICATION.} All tractable closed-form and data/arithmetic claims verify (C1, C3,
C4); the one DFT-energy-dependent claim (C2 Ti-Nb) is only partially reproduced because the
underlying VASP energies are out of scope and were surrogated (labelled). No fabrication: every
number is from paper text or a real script run.
\begin{center}
\textbf{Coverage: 6/10} \quad\quad \textbf{Agreement: 8/10}
\end{center}
Coverage is limited because the paper's core (DFT) is not independently recomputed. Agreement is
high on everything actually within scope.

\end{document}
